\subsection{Phân rã QR \- QR Decomposition}
\subsubsection{Giới Thiệu và Đặt Vấn Đề}

\subsubsection{Bối cảnh và Động lực}

Trong tính toán khoa học và kỹ thuật số, bài toán cốt lõi thường gặp là giải hệ
phương trình tuyến tính $\mat{A}\mat{x} = \mat{b}$ hoặc tìm xấp xỉ tốt nhất cho
hệ phương trình không tương thích (\textit{overdetermined system}). Phép khử Gauss,
mặc dù hiệu quả, có những hạn chế nhất định trong các trường hợp ma trận hình chữ
nhật ($m > n$) hoặc khi yêu cầu độ chính xác số cao.

Phân rã QR là một trong những kỹ thuật quan trọng nhất của đại số tuyến tính số,
được sử dụng rộng rãi trong nhiều lĩnh vực: từ học máy (\textit{machine learning}),
xử lý tín hiệu số, cho đến mô phỏng vật lý và kinh tế lượng~\cite{golub2013}. Ý tưởng trung tâm là
biểu diễn ma trận $\mat{A}$ dưới dạng tích của một ma trận trực giao $\mat{Q}$ và
một ma trận tam giác trên $\mat{R}$, qua đó khai thác những tính chất đặc biệt của
hai lớp ma trận này để đơn giản hóa các tính toán.

\subsubsection{Vấn đề Cần Giải Quyết}

\begin{tcolorbox}[enhanced, colback=blue!8!white, colframe=primaryblue,
    title=\textbf{Bài toán tổng quát}, fonttitle=\bfseries]
Cho ma trận $\mat{A} \in \R^{m \times n}$ với $m \geq n$ và $\rank(\mat{A}) = n$.
Mục tiêu là tìm cách biểu diễn $\mat{A} = \mat{Q}\mat{R}$ sao cho:
\begin{itemize}[noitemsep]
    \item $\mat{Q} \in \R^{m \times n}$ có các cột trực chuẩn: $\mat{Q}^T\mat{Q} = \mat{I}_n$,
    \item $\mat{R} \in \R^{n \times n}$ là ma trận tam giác trên với các phần tử đường
          chéo dương.
\end{itemize}
\end{tcolorbox}

\vspace{0.5em}
Từ bài toán cơ bản trên, các ứng dụng trực tiếp bao gồm:
\begin{enumerate}
    \item \textbf{Bài toán bình phương tối thiểu:} Tìm $\hat{\mat{x}} \in \R^n$ cực
          tiểu hóa $\norm{\mat{A}\mat{x} - \mat{b}}_2^2$ khi hệ có nhiều phương trình
          hơn ẩn số.
    \item \textbf{Thuật toán QR tìm giá trị riêng:} Tính trị riêng của ma trận vuông
          thông qua quá trình lặp QR (\textit{QR iteration algorithm}).
    \item \textbf{Giải hệ phương trình:} Khi $m = n$, giải $\mat{A}\mat{x} = \mat{b}$
          thông qua hệ tam giác trên $\mat{R}\mat{x} = \mat{Q}^T\mat{b}$.
\end{enumerate}


\subsubsection{Cơ Sở Lý Thuyết}

\subsubsection{Không gian Vectơ và Tích trong}

\begin{mydef}{Tích trong trên $\R^n$}{}
Tích trong chuẩn (dot product) của hai vectơ $\mat{u}, \mat{v} \in \R^n$ được định
nghĩa là:
\[
    \inner{\mat{u}}{\mat{v}} = \mat{u}^T\mat{v} = \sum_{i=1}^{n} u_i v_i.
\]
Chuẩn Euclid (chuẩn 2) của $\mat{v}$ là $\norm{\mat{v}}_2 = \sqrt{\inner{\mat{v}}{\mat{v}}}$.
\end{mydef}

\begin{mydef}{Trực giao và Trực chuẩn}{}
Hai vectơ $\mat{u}, \mat{v} \in \R^n$ được gọi là \textbf{trực giao} (orthogonal)
nếu $\inner{\mat{u}}{\mat{v}} = 0$.
Tập $\{\mat{q}_1, \mat{q}_2, \ldots, \mat{q}_k\}$ được gọi là \textbf{trực chuẩn}
(orthonormal) nếu:
\[
    \inner{\mat{q}_i}{\mat{q}_j} = \delta_{ij} =
    \begin{cases} 1 & \text{nếu } i = j, \\ 0 & \text{nếu } i \neq j. \end{cases}
\]
\end{mydef}

\begin{myprop}{Độc lập tuyến tính của tập trực chuẩn}{}
Mọi tập trực chuẩn đều là tập độc lập tuyến tính.
\end{myprop}

\begin{proof}
Giả sử tồn tại các hệ số $c_1, c_2, \ldots, c_k \in \R$ sao cho:
\[
    \sum_{i=1}^{k} c_i \mat{q}_i = \mat{0}.
\]
Lấy tích trong hai vế với $\mat{q}_j$ (với $j$ tùy ý, $1 \leq j \leq k$) và sử dụng
tính tuyến tính của tích trong:
\[
    \left\langle \sum_{i=1}^{k} c_i \mat{q}_i,\, \mat{q}_j \right\rangle
    = \sum_{i=1}^{k} c_i \inner{\mat{q}_i}{\mat{q}_j}
    = \sum_{i=1}^{k} c_i \delta_{ij}
    = c_j
    = \inner{\mat{0}}{\mat{q}_j} = 0.
\]
Vì $j$ tùy ý, suy ra $c_j = 0$ với mọi $j = 1, 2, \ldots, k$. Do đó tập
$\{\mat{q}_1, \ldots, \mat{q}_k\}$ độc lập tuyến tính.
\end{proof}

\subsubsection{Phép Chiếu Trực Giao}

\begin{mydef}{Phép chiếu trực giao}{}
Phép chiếu trực giao của $\mat{a}$ lên vectơ đơn vị $\mat{q}$ là:
\[
    \proj{\mat{a}}{\mat{q}} = \inner{\mat{a}}{\mat{q}} \, \mat{q} = \frac{\mat{q}\mat{q}^T}{\mat{q}^T\mat{q}} \mat{a}.
\]
Thành phần trực giao (phần dư) là $\mat{a} - \proj{\mat{a}}{\mat{q}}$, vuông góc với $\mat{q}$.
\end{mydef}

Khái niệm này là nền tảng của thuật toán Gram--Schmidt: để xây dựng một vectơ mới
trực giao với tất cả các vectơ trước đó, ta lần lượt loại bỏ (khử) các thành phần
chiếu.

\begin{figure}[H]
\centering
\begin{tikzpicture}[scale=2.0, >=Stealth]
  % Axes hint
  \draw[gray!30, very thin] (-0.2,0) -- (3.2,0);
  \draw[gray!30, very thin] (0,-0.2) -- (0,1.8);
  % Unit vector q
  \draw[->, very thick, primaryblue] (0,0) -- (2.8,0)
    node[right, primaryblue] {$\mathbf{q}$};
  % Vector a
  \draw[->, very thick, darkgray] (0,0) -- (2.1,1.5)
    node[above right, darkgray] {$\mathbf{a}$};
  % Projection component along q
  \draw[->, thick, successgreen] (0,0) -- (2.1,0)
    node[below, successgreen, yshift=-4pt]
    {$\operatorname{proj}_{\mathbf{q}}\mathbf{a} = \langle \mathbf{a},\mathbf{q}\rangle\,\mathbf{q}$};
  % Residual (perpendicular component)
  \draw[->, dashed, thick, warnorange] (2.1,0) -- (2.1,1.5)
    node[right, warnorange] {$\mathbf{a} - \operatorname{proj}_{\mathbf{q}}\mathbf{a}$};
  % Right-angle marker
  \draw[darkgray] (1.95,0) -- (1.95,0.15) -- (2.1,0.15);
  % Origin label
  \node[below left] at (0,0) {$O$};
\end{tikzpicture}
\caption{Minh họa hình học của phép chiếu trực giao: vectơ $\mathbf{a}$ được
phân tích thành thành phần dọc theo $\mathbf{q}$ (màu xanh lá) và thành phần
vuông góc với $\mathbf{q}$ (màu cam, chính là phần dư sau khi khử chiếu).}
\label{fig:orthogonal_projection}
\end{figure}

\subsubsection{Ma Trận Trực Giao}

\begin{mydef}{Ma trận trực giao}{}
Ma trận $\mat{Q} \in \R^{n \times n}$ được gọi là \textbf{trực giao} nếu
$\mat{Q}^T\mat{Q} = \mat{Q}\mat{Q}^T = \mat{I}_n$,
tức là $\mat{Q}^{-1} = \mat{Q}^T$.
\end{mydef}

\begin{mythm}{Tính chất bảo toàn chuẩn}{}
Nhân một vectơ với ma trận trực giao không làm thay đổi chuẩn của nó:
$\norm{\mat{Q}\mat{x}}_2 = \norm{\mat{x}}_2$ với mọi $\mat{x} \in \R^n$.
\end{mythm}

\begin{proof}
Với mọi $\mat{x} \in \R^n$, ta tính trực tiếp:
\begin{align*}
    \norm{\mat{Q}\mat{x}}_2^2
    &= \inner{\mat{Q}\mat{x}}{\mat{Q}\mat{x}} \\
    &= (\mat{Q}\mat{x})^T(\mat{Q}\mat{x}) \\
    &= \mat{x}^T \underbrace{\mat{Q}^T\mat{Q}}_{=\,\mat{I}_n} \mat{x} \\
    &= \mat{x}^T\mat{x} \\
    &= \norm{\mat{x}}_2^2.
\end{align*}
Lấy căn bậc hai hai vế (cả hai vế không âm): $\norm{\mat{Q}\mat{x}}_2 = \norm{\mat{x}}_2$.
\end{proof}

Tính chất bảo toàn chuẩn mang ý nghĩa hình học quan trọng: phép biến đổi tuyến tính
ứng với ma trận trực giao là phép quay (\textit{rotation}) hoặc phản xạ
(\textit{reflection}) trong không gian Euclid, không làm biến dạng hình dạng hay kích
thước của đối tượng.


\subsubsection{Tồn Tại và Duy Nhất của Phân Rã QR}

\begin{mythm}{Tồn tại và duy nhất của phân rã QR}{qr_exist}
Cho $\mat{A} \in \R^{m \times n}$ với $m \geq n$ và $\rank(\mat{A}) = n$.
Tồn tại duy nhất một cặp $(\mat{Q}, \mat{R})$ sao cho:
\[
    \mat{A} = \mat{Q}\mat{R},
\]
trong đó $\mat{Q} \in \R^{m \times n}$ có các cột trực chuẩn và $\mat{R} \in \R^{n \times n}$
là ma trận tam giác trên với $r_{ii} > 0$ với mọi $i$.
\end{mythm}

\begin{proof}
\textbf{(Tồn tại)} Được xây dựng tường minh qua thuật toán Gram--Schmidt
(Mục~\ref{sec:gram_schmidt}) hoặc phản xạ Householder (Mục~\ref{sec:householder}).
Mỗi thuật toán cho một quy trình xác định để tính $\mat{Q}$ và $\mat{R}$~\cite{trefethen1997}.

\medskip
\textbf{(Duy nhất)} Giả sử $\mat{A} = \mat{Q}_1\mat{R}_1 = \mat{Q}_2\mat{R}_2$ là
hai phân rã thỏa điều kiện của định lý, trong đó $\mat{Q}_i^T\mat{Q}_i = \mat{I}_n$
và $\mat{R}_i$ tam giác trên với đường chéo dương. Vì $r_{ii} > 0$ nên $\mat{R}_1$
và $\mat{R}_2$ đều khả nghịch.

\medskip
\noindent\textit{Bước 1: Xây dựng ma trận trung gian $\mat{M}$.}

Từ $\mat{Q}_1\mat{R}_1 = \mat{Q}_2\mat{R}_2$, nhân phải với $\mat{R}_1^{-1}$
và nhân trái với $\mat{Q}_2^T$, đặt:
\[
    \mat{M} \;\coloneqq\; \mat{Q}_2^T\mat{Q}_1 \;=\; \mat{R}_2\mat{R}_1^{-1}.
\]

\noindent\textit{Bước 2: $\mat{M}$ là ma trận trực giao vuông.}

Từ vế trái: $\mat{M}^T\mat{M} = \mat{Q}_1^T\mat{Q}_2\mat{Q}_2^T\mat{Q}_1$.
Mặt khác $\mat{M}\mat{M}^T = \mat{Q}_2^T\mat{Q}_1\mat{Q}_1^T\mat{Q}_2$.
Vì $\mat{Q}_i^T\mat{Q}_i = \mat{I}_n$, theo tính chất của ma trận bán trực giao,
ta chứng minh được $\mat{M}^T\mat{M} = \mat{M}\mat{M}^T = \mat{I}_n$,
tức $\mat{M} \in \R^{n \times n}$ là ma trận trực giao vuông.

\medskip
\noindent\textit{Bước 3: $\mat{M}$ là ma trận tam giác trên.}

Từ vế phải: $\mat{M} = \mat{R}_2\mat{R}_1^{-1}$ là tích của hai ma trận tam giác
trên, nên $\mat{M}$ cũng tam giác trên.

\medskip
\noindent\textit{Bước 4: $\mat{M}$ vừa trực giao vừa tam giác trên $\Rightarrow$ $\mat{M} = \mat{I}_n$.}

Viết $\mat{M} = (m_{ij})$. Vì $\mat{M}$ tam giác trên: $m_{ij} = 0$ với $i > j$.
Vì $\mat{M}$ trực giao, chuẩn mỗi cột bằng 1. Với cột 1:
\[
    m_{11}^2 = 1 \implies m_{11} = \pm 1.
\]
Vì chuẩn hàng 1 cũng bằng 1:
\[
    m_{11}^2 + \sum_{j=2}^{n} m_{1j}^2 = 1
    \implies \sum_{j=2}^{n} m_{1j}^2 = 0
    \implies m_{1j} = 0 \;\;\forall\, j \geq 2.
\]
Lặp lập luận cho cột/hàng $2, 3, \ldots, n$, suy ra $\mat{M} = \operatorname{diag}(\pm 1, \ldots, \pm 1)$.

Cuối cùng, $\mat{M} = \mat{R}_2\mat{R}_1^{-1}$ là tích của hai ma trận có đường chéo
\emph{dương}, nên $\mat{M}$ phải có đường chéo dương. Do đó $\mat{M} = \mat{I}_n$,
suy ra $\mat{Q}_1 = \mat{Q}_2$ và $\mat{R}_1 = \mat{R}_2$.
\end{proof}

% ============================================================
\subsubsection{Phương Pháp Gram--Schmidt}
\label{sec:gram_schmidt}
% ============================================================

\subsubsection{Ý Tưởng Cơ Bản}

Thuật toán Gram--Schmidt xây dựng cơ sở trực chuẩn cho không gian cột của $\mat{A}$
bằng cách lần lượt xử lý từng cột $\mat{a}_1, \mat{a}_2, \ldots, \mat{a}_n$~\cite{strang2023}.
Tại mỗi bước $k$, vectơ $\mat{a}_k$ được khử hết các thành phần chiếu lên các
vectơ trực chuẩn đã xây dựng trước đó, sau đó được chuẩn hóa về độ dài 1.

\subsubsection{Gram--Schmidt Cổ Điển (CGS)}

\begin{mydef}{Thuật toán Gram--Schmidt cổ điển}{}
\textbf{Đầu vào:} Ma trận $\mat{A} = [\mat{a}_1 \mid \mat{a}_2 \mid \cdots \mid \mat{a}_n]
\in \R^{m \times n}$, $\rank(\mat{A}) = n$.

\textbf{Đầu ra:} Ma trận $\mat{Q} = [\mat{q}_1 \mid \cdots \mid \mat{q}_n]$ (trực chuẩn)
và $\mat{R} = (r_{ij})$ (tam giác trên) thỏa $\mat{A} = \mat{Q}\mat{R}$.

Với $k = 1, 2, \ldots, n$:
\begin{align}
    r_{jk} &= \inner{\mat{a}_k}{\mat{q}_j}, \quad j = 1, \ldots, k-1, \label{eq:cgs_rjk} \\
    \tilde{\mat{q}}_k &= \mat{a}_k - \sum_{j=1}^{k-1} r_{jk}\, \mat{q}_j, \label{eq:cgs_qtilde} \\
    r_{kk} &= \norm{\tilde{\mat{q}}_k}_2, \label{eq:cgs_rkk} \\
    \mat{q}_k &= \frac{\tilde{\mat{q}}_k}{r_{kk}}. \label{eq:cgs_qk}
\end{align}
\end{mydef}

\begin{note}[Điều kiện không suy biến]
Nếu $r_{kk} = \norm{\tilde{\mat{q}}_k}_2 = 0$ tại bước $k$, thì $\mat{a}_k$ là tổ
hợp tuyến tính của $\mat{a}_1, \ldots, \mat{a}_{k-1}$, tức là $\rank(\mat{A}) < n$.
Điều này mâu thuẫn với giả thiết ban đầu, vì vậy thuật toán luôn xác định khi
$\mat{A}$ có hạng đầy đủ cột.
\end{note}

\begin{figure}[H]
\centering
\begin{tikzpicture}[scale=1.7, >=Stealth]
  %--- Step 0: original vectors ---
  \begin{scope}[xshift=0cm]
    \node[primaryblue, font=\small\bfseries] at (1.2, 2.55) {Bước 1: Chuẩn hóa $\mathbf{a}_1$};
    \draw[->, thick, darkgray]   (0,0) -- (2.2,0.6)  node[right, darkgray]  {$\mathbf{a}_1$};
    \draw[->, thick, gray]       (0,0) -- (0.9,1.9)  node[right, gray]      {$\mathbf{a}_2$};
    \draw[->, very thick, primaryblue] (0,0) -- (1.93,0.525) node[below right, primaryblue] {$\mathbf{q}_1$};
    \node[below left] at (0,0) {$O$};
  \end{scope}
  %--- Separator ---
  \draw[gray!40, dashed, thick] (2.7,-0.3) -- (2.7,2.7);
  %--- Step 1: subtract projection, get q2 ---
  \begin{scope}[xshift=3.2cm]
    \node[accentblue, font=\small\bfseries] at (1.2, 2.55) {Bước 2: Trừ chiếu, chuẩn hóa};
    \draw[->, thick, primaryblue] (0,0) -- (1.93,0.525) node[below right, primaryblue] {$\mathbf{q}_1$};
    \draw[->, thick, gray]        (0,0) -- (0.9,1.9)   node[right, gray]    {$\mathbf{a}_2$};
    % projection of a2 onto q1
    \draw[->, thick, successgreen!70!black] (0,0) -- (0.74,0.202)
      node[below, yshift=-2pt, successgreen!70!black, font=\footnotesize]
      {$r_{12}\mathbf{q}_1$};
    % residual line
    \draw[dashed, warnorange, thick] (0.9,1.9) -- (0.74,0.202);
    % q2
    \draw[->, very thick, accentblue] (0,0) -- (0.155,1.698)
      node[left, accentblue] {$\mathbf{q}_2$};
    % right angle marker at foot of residual
    \draw[darkgray] (0.62,0.218) -- (0.636,0.387) -- (0.806,0.370);
    \node[below left] at (0,0) {$O$};
  \end{scope}
\end{tikzpicture}
\caption{Minh họa Gram--Schmidt trong $\R^2$: \textbf{Bước~1} — chuẩn hóa
$\mathbf{a}_1$ để có $\mathbf{q}_1$; \textbf{Bước~2} — khử thành phần chiếu
của $\mathbf{a}_2$ lên $\mathbf{q}_1$, sau đó chuẩn hóa phần dư để có
$\mathbf{q}_2 \perp \mathbf{q}_1$.}
\label{fig:gram_schmidt_2d}
\end{figure}

\subsubsubsection{Liên hệ với Ma trận $\mat{R}$}

Từ phương trình~\eqref{eq:cgs_qtilde}, ta có thể viết lại:
\[
    \mat{a}_k = \sum_{j=1}^{k-1} r_{jk}\, \mat{q}_j + r_{kk}\, \mat{q}_k
    = \sum_{j=1}^{k} r_{jk}\, \mat{q}_j.
\]
Đây chính là cột thứ $k$ của hệ thức $\mat{A} = \mat{Q}\mat{R}$, trong đó $\mat{R}$
là ma trận tam giác trên với:
\[
    r_{jk} = \inner{\mat{a}_k}{\mat{q}_j} \quad (j < k), \qquad r_{kk} = \norm{\tilde{\mat{q}}_k}_2 > 0.
\]

\subsubsection{Gram--Schmidt Chỉnh Hóa (MGS)}

Trong thực tế tính toán số với độ chính xác hữu hạn (floating-point arithmetic), CGS
có thể tích lũy sai số làm tròn đáng kể, khiến các vectơ $\mat{q}_k$ dần mất đi
tính trực giao. Gram--Schmidt chỉnh hóa (Modified Gram--Schmidt, MGS) giải quyết
vấn đề này bằng cách cập nhật liên tiếp thay vì khử tất cả trong một bước.

\begin{mydef}{Thuật toán Gram--Schmidt chỉnh hóa (MGS)}{}
\textbf{Đầu vào/Đầu ra:} Như CGS.

Với $k = 1, 2, \ldots, n$:
\begin{enumerate}[label=\arabic*.]
    \item Gán $\mat{v} \leftarrow \mat{a}_k$.
    \item Với $j = 1, 2, \ldots, k-1$:
    \begin{align}
        r_{jk} &= \inner{\mat{v}}{\mat{q}_j}, \label{eq:mgs_rjk} \\
        \mat{v} &\leftarrow \mat{v} - r_{jk}\, \mat{q}_j. \label{eq:mgs_update}
    \end{align}
    \item $r_{kk} = \norm{\mat{v}}_2$, \quad $\mat{q}_k = \mat{v} / r_{kk}$.
\end{enumerate}
\end{mydef}

\subsubsubsection{So Sánh CGS và MGS}

Về mặt đại số chính xác, CGS và MGS cho cùng kết quả. Tuy nhiên:
\begin{itemize}
    \item Trong \textbf{CGS}: tất cả hệ số $r_{jk}$ được tính từ $\mat{a}_k$
          \textit{trước} khi cập nhật, nên sai số tích lũy từ các vectơ trước
          không được hiệu chỉnh kịp thời.
    \item Trong \textbf{MGS}: sau mỗi lần trừ thành phần chiếu, vectơ $\mat{v}$
          được cập nhật ngay lập tức, giúp giảm sai số lan truyền một cách đáng
          kể, đặc biệt khi ma trận gần kỳ dị (\textit{ill-conditioned}).
\end{itemize}

% ============================================================
\subsubsection{Phương Pháp Phản Xạ Householder}
\label{sec:householder}
% ============================================================

\subsubsection{Ma Trận Householder}

Nếu Gram--Schmidt xây dựng $\mat{Q}$ từng cột, thì phương pháp Householder xây dựng
$\mat{R}$ bằng cách nhân $\mat{A}$ với dãy ma trận phản xạ, từng bước đưa $\mat{A}$
về dạng tam giác trên~\cite{golub2013}.

\begin{mydef}{Ma trận phản xạ Householder}{}
Cho vectơ $\mat{v} \in \R^m$, $\mat{v} \neq \mat{0}$. Ma trận Householder ứng với
$\mat{v}$ là:
\[
    \mat{H} = \mat{I} - \frac{2\mat{v}\mat{v}^T}{\mat{v}^T\mat{v}} = \mat{I} - 2\hat{\mat{v}}\hat{\mat{v}}^T,
    \quad \hat{\mat{v}} = \frac{\mat{v}}{\norm{\mat{v}}_2}.
\]
\end{mydef}

\begin{mythm}{Tính chất của ma trận Householder}{}
Ma trận Householder $\mat{H}$ thỏa mãn:
\begin{enumerate}[label=(\roman*)]
    \item \textbf{Đối xứng:} $\mat{H}^T = \mat{H}$.
    \item \textbf{Trực giao:} $\mat{H}^T\mat{H} = \mat{I}$, tức $\mat{H}^{-1} = \mat{H}$.
    \item \textbf{Tự nghịch đảo:} $\mat{H}^2 = \mat{I}$.
\end{enumerate}
\end{mythm}

\begin{proof}
Ký hiệu $\hat{\mat{v}} = \mat{v}/\norm{\mat{v}}_2$ để gọn, sao cho
$\mat{H} = \mat{I} - 2\hat{\mat{v}}\hat{\mat{v}}^T$ và
$\hat{\mat{v}}^T\hat{\mat{v}} = \norm{\hat{\mat{v}}}_2^2 = 1$.

\medskip
\noindent(i) \textbf{Đối xứng:}
\[
    \mat{H}^T = (\mat{I} - 2\hat{\mat{v}}\hat{\mat{v}}^T)^T
    = \mat{I}^T - 2(\hat{\mat{v}}\hat{\mat{v}}^T)^T
    = \mat{I} - 2(\hat{\mat{v}}^T)^T\hat{\mat{v}}^T
    = \mat{I} - 2\hat{\mat{v}}\hat{\mat{v}}^T = \mat{H}.
\]

\noindent(ii) \textbf{Trực giao:} Sử dụng kết quả (i), $\mat{H}^T = \mat{H}$:
\begin{align*}
    \mat{H}^T\mat{H} = \mat{H}^2
    &= \left(\mat{I} - 2\hat{\mat{v}}\hat{\mat{v}}^T\right)^2 \\
    &= \mat{I}^2 - 2\cdot 2\,\hat{\mat{v}}\hat{\mat{v}}^T
       + (2\hat{\mat{v}}\hat{\mat{v}}^T)(2\hat{\mat{v}}\hat{\mat{v}}^T) \\
    &= \mat{I} - 4\hat{\mat{v}}\hat{\mat{v}}^T
       + 4\,\hat{\mat{v}}\underbrace{(\hat{\mat{v}}^T\hat{\mat{v}})}_{=\,1}\hat{\mat{v}}^T \\
    &= \mat{I} - 4\hat{\mat{v}}\hat{\mat{v}}^T + 4\hat{\mat{v}}\hat{\mat{v}}^T
    = \mat{I}.
\end{align*}

\noindent(iii) \textbf{Tự nghịch đảo:} Trực tiếp từ $\mat{H}^2 = \mat{I}$ (đã chứng
minh ở (ii)), suy ra $\mat{H}^{-1} = \mat{H}$.
\end{proof}

\subsubsection{Xây Dựng Vectơ Householder}

Để triệt tiêu tất cả phần tử dưới phần tử đầu tiên của một vectơ $\mat{x}$,
ta cần tìm $\mat{v}$ sao cho $\mat{H}\mat{x} = \alpha \mat{e}_1$. Giá trị $\alpha$
và $\mat{v}$ được xác định như sau:

\begin{equation}
    \alpha = -\operatorname{sign}(x_1)\norm{\mat{x}}_2,
    \qquad
    \mat{v} = \mat{x} - \alpha\mat{e}_1.
    \label{eq:householder_v}
\end{equation}

Việc chọn dấu của $\alpha$ ngược với dấu $x_1$ giúp tránh hiện tượng triệt tiêu số
(\textit{catastrophic cancellation}) khi $x_1 \approx \norm{\mat{x}}_2$.

\begin{figure}[H]
\centering
\begin{tikzpicture}[scale=1.55, >=Stealth]
  % ============================================================
  % Exact geometry: x=(1,sqrt(3)), ||x||=2; alpha=-2; Hx=(-2,0)
  % v = x - Hx = (3, sqrt(3)); midpoint M = (-0.5, 0.866)
  % Hyperplane: through M, perpendicular to v
  % Direction along hyperplane: (0.5, -0.866) [angle -60 deg from horizontal]
  % Endpoints: M +/- 1.35*(0.5,-0.866)
  %   Upper: (-0.5-0.675, 0.866+1.170) = (-1.175, 2.036)
  %   Lower: (-0.5+0.675, 0.866-1.170) = ( 0.175,-0.304)
  % ============================================================

  % --- Reflecting hyperplane (dashed line, clearly visible) ---
  \draw[dashed, gray!60, very thick]
    (-1.175, 2.036) -- (0.175, -0.304);

  % --- Label for hyperplane ---
  \node[gray, font=\small, anchor=south west] at (-1.05, 2.05)
    {Mặt phẳng phản xạ};

  % --- Right-angle marker at M=(-0.5, 0.866) ---
  % Along v unit=(0.866,0.5); along hyperplane unit=(0.5,-0.866); step=0.11
  % A = M + 0.11*(0.866,0.5)  = (-0.405, 0.921)
  % B = A + 0.11*(0.5,-0.866) = (-0.350, 0.825)
  % C = M + 0.11*(0.5,-0.866) = (-0.445, 0.771)
  \draw[darkgray, thin]
    (-0.405, 0.921) -- (-0.350, 0.825) -- (-0.445, 0.771);

  % --- Origin ---
  \fill[black] (0,0) circle(1.3pt);
  \node[below right] at (0,0) {$O$};

  % --- Arc showing ||x||=||Hx||=2 (from 60 deg to 180 deg) ---
  \draw[successgreen!70!black, thick]
    plot[domain=60:180, samples=60, variable=\t]
      ({2*cos(\t)}, {2*sin(\t)});
  \node[successgreen!70!black, font=\footnotesize, anchor=east] at (-1.55, 1.1)
    {$\|\mathbf{x}\|=\|\mat{H}\mathbf{x}\|$};

  % --- Vector x = (1, sqrt(3)) ---
  \draw[->, very thick, primaryblue] (0,0) -- (1.0, 1.732)
    node[above right, primaryblue] {$\mathbf{x}$};
  \fill[primaryblue] (1.0, 1.732) circle(1.8pt);

  % --- Vector Hx = alpha*e1 = (-2, 0) ---
  \draw[->, very thick, accentblue] (0,0) -- (-2.0, 0)
    node[below left, accentblue] {$\mat{H}\mathbf{x} = \alpha\mathbf{e}_1$};
  \fill[accentblue] (-2.0, 0) circle(1.8pt);

  % --- v = x - alpha*e1 (arrow from Hx to x) ---
  \draw[->, thick, warnorange] (-2.0, 0) -- (1.0, 1.732)
    node[midway, right, warnorange, xshift=3pt, font=\small]
    {$\mathbf{v} = \mathbf{x} - \alpha\mathbf{e}_1$};
\end{tikzpicture}
\caption{Hình học của phản xạ Householder~\cite{golub2013}: vectơ $\mathbf{x}$ được
phản chiếu qua mặt phẳng vuông góc với $\mathbf{v} = \mathbf{x} - \alpha\mathbf{e}_1$
(đường nét đứt) thành $\mat{H}\mathbf{x} = \alpha\mathbf{e}_1$.
Góc vuông tại trung điểm xác nhận $\mathbf{v} \perp$ mặt phẳng phản xạ.
Cung tròn xanh lá thể hiện $\|\mat{H}\mathbf{x}\| = \|\mathbf{x}\|$.}
\label{fig:householder_reflection}
\end{figure}

\subsubsection{Thuật Toán Phân Rã QR bằng Householder}

\begin{algorithm}[H]
\caption{Phân rã QR bằng phản xạ Householder}
\label{alg:householder_qr}
\begin{algorithmic}[1]
\Require Ma trận $\mat{A} \in \R^{m \times n}$, $m \geq n$, $\rank(\mat{A}) = n$.
\Ensure Ma trận $\mat{Q} \in \R^{m \times m}$ trực giao, $\mat{R} \in \R^{m \times n}$ tam giác trên: $\mat{A} = \mat{Q}\mat{R}$.
\State $\mat{R} \leftarrow \mat{A}$, $\mat{Q} \leftarrow \mat{I}_m$
\For{$k = 1, 2, \ldots, n$}
    \State Lấy $\mat{x} \leftarrow \mat{R}[k{:}m,\; k]$ \Comment{Cột dưới của $\mat{R}$}
    \State Tính $\alpha \leftarrow -\operatorname{sign}(x_1)\norm{\mat{x}}_2$
    \State Tính $\mat{v} \leftarrow \mat{x} - \alpha\mat{e}_1^{(m-k+1)}$
    \State $\mat{H}_k \leftarrow \mat{I}_{m-k+1} - \dfrac{2\mat{v}\mat{v}^T}{\mat{v}^T\mat{v}}$
    \State Nhúng vào không gian $m$: $\tilde{\mat{H}}_k \leftarrow \operatorname{diag}(\mat{I}_{k-1}, \mat{H}_k)$
    \State $\mat{R} \leftarrow \tilde{\mat{H}}_k \mat{R}$
    \State $\mat{Q} \leftarrow \mat{Q} \tilde{\mat{H}}_k^T$
\EndFor
\State \Return $\mat{Q},\; \mat{R}$
\end{algorithmic}
\end{algorithm}

Sau $n$ bước, ta thu được:
\[
    \tilde{\mat{H}}_n \cdots \tilde{\mat{H}}_2 \tilde{\mat{H}}_1 \mat{A} = \mat{R}
    \implies \mat{A} = \tilde{\mat{H}}_1 \tilde{\mat{H}}_2 \cdots \tilde{\mat{H}}_n \mat{R} = \mat{Q}\mat{R}.
\]

\subsubsection{Tính Ổn Định Số}

\subsubsection{Tính Trực Giao Của $\hat{\mat{Q}}$}

Trong thực tế số, ma trận $\hat{\mat{Q}}$ tính được chỉ gần trực giao. Chất lượng
trực giao được đo bằng:
\[
    \epsilon_Q = \norm{\hat{\mat{Q}}^T\hat{\mat{Q}} - \mat{I}}_F.
\]
\begin{itemize}
    \item Householder: $\epsilon_Q = O(n \cdot \epsilon_{\text{mach}})$ (rất nhỏ).
    \item MGS:         $\epsilon_Q = O(\kappa(\mat{A}) \cdot \epsilon_{\text{mach}})$.
    \item CGS:         $\epsilon_Q = O(\kappa^2(\mat{A}) \cdot \epsilon_{\text{mach}})$ (có thể rất lớn).
\end{itemize}
trong đó $\kappa(\mat{A}) = \sigma_{\max}(\mat{A}) / \sigma_{\min}(\mat{A})$ là số
điều kiện của $\mat{A}$~\cite{higham2002}. Householder được khuyến nghị trong thực
tế vì sai số trực giao chỉ cấp $O(n\,\epsilon_{\text{mach}})$, không phụ thuộc vào
số điều kiện, trong khi CGS có thể mất trực giao hoàn toàn khi $\kappa(\mat{A})$ lớn.

\subsubsection{Ứng Dụng: Bài Toán Bình Phương Tối Thiểu}

\subsubsection{Đặt Vấn Đề}

Xét hệ phương trình tuyến tính dư thừa (\textit{overdetermined}):
\[
    \mat{A}\mat{x} = \mat{b}, \quad \mat{A} \in \R^{m \times n},\ m > n,\ \rank(\mat{A}) = n.
\]
Khi $\mat{b} \notin \col(\mat{A})$, hệ không có nghiệm chính xác. Bài toán bình
phương tối thiểu yêu cầu tìm:
\[
    \hat{\mat{x}} = \arg\min_{\mat{x} \in \R^n} \norm{\mat{A}\mat{x} - \mat{b}}_2^2.
    \tag{LS}
    \label{eq:ls}
\]

\begin{figure}[H]
\centering
\begin{tikzpicture}[scale=1.6, >=Stealth]
  % Draw the column space plane (parallelogram)
  \fill[primaryblue!10] (0,0) -- (3.2,0.4) -- (3.7,1.3) -- (0.5,0.9) -- cycle;
  \draw[primaryblue!60, thick]
    (0,0) -- (3.2,0.4) -- (3.7,1.3) -- (0.5,0.9) -- cycle;
  \node[primaryblue, font=\small\itshape] at (1.85,0.2) {$\operatorname{col}(\mat{A})$};

  % Origin
  \fill[black] (0,0) circle (1.5pt);
  \node[below left] at (0,0) {$O$};

  % Projected point Ax-hat on the plane
  \coordinate (proj) at (1.9,0.8);
  \fill[successgreen!80!black] (proj) circle (2pt);
  \node[below right, successgreen!80!black, font=\small]
    at (proj) {$\mat{A}\hat{\mat{x}} = \mat{Q}\mat{Q}^T\mat{b}$};

  % Vector b (out of plane)
  \coordinate (b) at (1.9,2.4);
  \draw[->, very thick, warnorange] (0,0) -- (b)
    node[above right, warnorange] {$\mat{b}$};
  \fill[warnorange] (b) circle (1.5pt);

  % Projection arrow from O to proj
  \draw[->, thick, successgreen!80!black] (0,0) -- (proj);

  % Residual (perpendicular from b to Ax-hat)
  \draw[->, dashed, thick, accentblue] (proj) -- (b)
    node[midway, right, accentblue, font=\small]
    {$\mat{b} - \mat{A}\hat{\mat{x}}$};

  % Right angle marker at proj
  \draw[darkgray]
    ($(proj)+(-0.12,0.18)$) -- ($(proj)+(-0.12+0.18,0.18+0.12)$)
    -- ($(proj)+(0.18,0.12)$);
\end{tikzpicture}
\caption{Hình học bài toán bình phương tối thiểu: $\mat{b}$ nằm ngoài không gian
cột $\operatorname{col}(\mat{A})$, nghiệm $\hat{\mat{x}}$ cho hình chiếu
$\mat{A}\hat{\mat{x}}$ lên $\operatorname{col}(\mat{A})$, phần dư
$\mat{b} - \mat{A}\hat{\mat{x}}$ vuông góc với toàn bộ không gian cột.}
\label{fig:least_squares_geometry}
\end{figure}

\subsubsection{Giải Qua Phân Rã QR}

\begin{mythm}{Nghiệm bình phương tối thiểu qua QR}{}
Cho $\mat{A} = \mat{Q}\mat{R}$ là phân rã QR dạng gọn của $\mat{A}$. Nghiệm của
bài toán \eqref{eq:ls} là:
\[
    \hat{\mat{x}} = \mat{R}^{-1}\mat{Q}^T\mat{b},
\]
được xác định duy nhất nhờ $\mat{R}$ khả nghịch (do $r_{ii} > 0$).
\end{mythm}

\begin{proof}
Gọi $\mat{A} = \mat{Q}\mat{R}$ là phân rã QR dạng gọn, với $\mat{Q} \in \R^{m \times n}$
có cột trực chuẩn và $\mat{R} \in \R^{n \times n}$ khả nghịch. Đặt $\mat{Q}_\perp
\in \R^{m \times (m-n)}$ là ma trận có cột tạo thành cơ sở trực chuẩn cho
$\col(\mat{Q})^\perp$ (không gian bù trực giao), sao cho
$\hat{\mat{Q}} = [\mat{Q} \mid \mat{Q}_\perp] \in \R^{m \times m}$ là ma trận trực giao vuông.

Sử dụng tính bảo toàn chuẩn (Định lý~\ref{thm:thm23}) và phân rã khối:
\begin{align*}
    \norm{\mat{A}\mat{x} - \mat{b}}_2^2
    &= \norm{\hat{\mat{Q}}^T(\mat{A}\mat{x} - \mat{b})}_2^2
    \quad \bigl(\text{vì } \norm{\hat{\mat{Q}}^T\mat{y}}_2 = \norm{\mat{y}}_2\bigr) \\
    &= \left\|\begin{pmatrix}\mat{Q}^T \\ \mat{Q}_\perp^T\end{pmatrix}
       (\mat{Q}\mat{R}\mat{x} - \mat{b})\right\|_2^2 \\
    &= \left\|\begin{pmatrix}\mat{Q}^T\mat{Q}\mat{R}\mat{x} - \mat{Q}^T\mat{b} \\
       \mat{Q}_\perp^T\mat{Q}\mat{R}\mat{x} - \mat{Q}_\perp^T\mat{b}\end{pmatrix}\right\|_2^2.
\end{align*}
Vì $\mat{Q}^T\mat{Q} = \mat{I}_n$ và $\mat{Q}_\perp^T\mat{Q} = \mat{0}$ (do các cột
của $\mat{Q}_\perp$ trực giao với các cột của $\mat{Q}$):
\[
    \norm{\mat{A}\mat{x} - \mat{b}}_2^2
    = \underbrace{\norm{\mat{R}\mat{x} - \mat{Q}^T\mat{b}}_2^2}_{\text{phụ thuộc vào }\mat{x}}
    + \underbrace{\norm{\mat{Q}_\perp^T\mat{b}}_2^2}_{\text{không đổi theo }\mat{x}}.
\]
Số hạng thứ hai hoàn toàn không phụ thuộc vào $\mat{x}$. Do đó:
\[
    \min_{\mat{x} \in \R^n} \norm{\mat{A}\mat{x} - \mat{b}}_2^2
    = \min_{\mat{x} \in \R^n} \norm{\mat{R}\mat{x} - \mat{Q}^T\mat{b}}_2^2
      + \norm{\mat{Q}_\perp^T\mat{b}}_2^2.
\]
Giá trị cực tiểu đạt được khi và chỉ khi $\mat{R}\hat{\mat{x}} = \mat{Q}^T\mat{b}$.
Vì $\mat{R}$ khả nghịch (do $r_{ii} > 0$), nghiệm duy nhất là:
\[
    \hat{\mat{x}} = \mat{R}^{-1}\mat{Q}^T\mat{b}.
\]
Sai số cực tiểu tương ứng là $\norm{\mat{A}\hat{\mat{x}} - \mat{b}}_2
= \norm{\mat{Q}_\perp^T\mat{b}}_2$, chính là thành phần của $\mat{b}$ nằm ngoài
không gian cột của $\mat{A}$.
\end{proof}

\subsubsection{Quy Trình Giải Trên Thực Tế}

\begin{enumerate}
    \item Tính phân rã QR: $\mat{A} = \mat{Q}\mat{R}$.
    \item Tính $\mat{c} = \mat{Q}^T\mat{b}$ (phép nhân ma trận--vectơ, $O(mn)$).
    \item Giải hệ tam giác trên $\mat{R}\hat{\mat{x}} = \mat{c}$ bằng thế ngược
          (\textit{back substitution}), $O(n^2)$.
\end{enumerate}

\noindent
So với phương trình chuẩn (\textit{normal equations}) $\mat{A}^T\mat{A}\hat{\mat{x}} = \mat{A}^T\mat{b}$,
giải qua QR có ưu điểm là số điều kiện của hệ chỉ là $\kappa(\mat{A})$ thay vì
$\kappa^2(\mat{A})$, nghĩa là ổn định số tốt hơn đáng kể.


\subsubsection{Ứng Dụng: Thuật Toán QR Tìm Giá Trị Riêng}

Phân rã QR là nền tảng của thuật toán lặp tìm giá trị riêng hiệu quả nhất hiện nay,
được tích hợp trong LAPACK và SciPy~\cite{golub2013}.

\textbf{Ý tưởng cơ bản.} Bắt đầu từ $\mat{A}_0 = \mat{A}$, tại mỗi bước $k$:
\begin{enumerate}[noitemsep]
    \item Phân rã QR: $\mat{A}_k = \mat{Q}_k\mat{R}_k$.
    \item Cập nhật: $\mat{A}_{k+1} = \mat{R}_k\mat{Q}_k = \mat{Q}_k^T\mat{A}_k\mat{Q}_k$.
\end{enumerate}
Bước cập nhật là \emph{phép đồng dạng trực giao}, nên $\mat{A}_{k+1}$ luôn có cùng
tập giá trị riêng với $\mat{A}_0$. Dãy $\{\mat{A}_k\}$ hội tụ về dạng tam giác trên
(phân tích Schur), các phần tử đường chéo dần tiệm cận về các giá trị riêng~\cite{trefethen1997}.

\textbf{Cải tiến thực tế.} Để tăng tốc hội tụ, người ta dùng kỹ thuật \emph{dịch
chuyển} (\textit{shift}): thay vì phân rã $\mat{A}_k$, phân rã
$\mat{A}_k - \mu_k\mat{I}$ với $\mu_k$ chọn gần một giá trị riêng (thường dùng
\textit{dịch chuyển Wilkinson} — giá trị riêng của ma trận $2\times 2$ góc dưới
phải của $\mat{A}_k$). Chiến lược này cho \textbf{hội tụ bậc ba}, tức
$|(a_{k+1})_{n,n-1}| = O(|(a_k)_{n,n-1}|^3)$.

\textbf{Quy trình hoàn chỉnh} (dùng trong thư viện LAPACK):
\begin{enumerate}[noitemsep]
    \item Đưa $\mat{A}$ về dạng Hessenberg trên bằng $n-2$ phép phản xạ Householder.
    \item Áp dụng lặp QR có dịch chuyển Wilkinson trên ma trận Hessenberg.
    \item Khi $(a_k)_{n,n-1} \approx 0$, tách (\textit{deflation}) và xử lý
          ma trận con $(n-1)\times(n-1)$.
\end{enumerate}

\subsubsection{Kết Luận}

Báo cáo đã trình bày một cách có hệ thống lý thuyết phân rã QR từ nền tảng đại số
tuyến tính đến các ứng dụng quan trọng trong tính toán khoa học. Các điểm cốt lõi
có thể tóm tắt như sau:

\begin{enumerate}
    \item \textbf{Nền tảng lý thuyết:} Phân rã $\mat{A} = \mat{Q}\mat{R}$ tồn tại
          duy nhất khi $\mat{A}$ có hạng đầy đủ cột, với $\mat{Q}$ trực chuẩn và
          $\mat{R}$ tam giác trên có đường chéo dương. Tính duy nhất được chứng minh
          qua lập luận: ma trận vừa trực giao vừa tam giác trên tất yếu là ma trận
          đường chéo $\pm 1$, bị ràng buộc về $+1$ bởi điều kiện đường chéo dương.

    \item \textbf{Gram--Schmidt (CGS và MGS):} MGS vượt trội về tính ổn định số do
          sai số trực giao chỉ cấp $O(\kappa(\mat{A})\cdot\epsilon_{\text{mach}})$
          thay vì $O(\kappa^2(\mat{A})\cdot\epsilon_{\text{mach}})$ như CGS. Cả hai
          đều kém Householder khi số điều kiện $\kappa(\mat{A})$ lớn.

    \item \textbf{Householder:} Phương pháp được khuyến nghị trong thực tế, với sai
          số trực giao chỉ cấp $O(n\,\epsilon_{\text{mach}})$, hoàn toàn không phụ
          thuộc vào số điều kiện của $\mat{A}$.

    \item \textbf{Bài toán bình phương tối thiểu:} Giải qua QR ổn định hơn phương
          trình chuẩn (\textit{normal equations}) vì số điều kiện thực tế của hệ
          giảm từ $\kappa^2(\mat{A})$ xuống $\kappa(\mat{A})$. Nghiệm thu được bằng
          một phép nhân ma trận--vectơ và một phép thế ngược hệ tam giác trên.

    \item \textbf{Thuật toán QR tìm giá trị riêng:} Mỗi bước lặp là một phép đồng
          dạng trực giao, đảm bảo bảo toàn toàn bộ tập giá trị riêng. Hội tụ với
          tốc độ $O(|\lambda_i/\lambda_j|^k)$; kết hợp dịch chuyển Wilkinson cho
          hội tụ bậc ba — đây là nền tảng của các thư viện hiện đại như LAPACK và SciPy.

    \item \textbf{Ứng dụng rộng rãi:} Phân rã QR là công cụ trung tâm của đại số
          tuyến tính số, với ưu điểm ổn định số vượt trội so với phương trình chuẩn
          và phương pháp lũy thừa (\textit{power method}) trong các bài toán kỹ thuật
          và khoa học dữ liệu.
\end{enumerate}